\chapter{Conclusions}

\section*{Future Reccomendations}

The results indicated that the design was effective in providing reliable monitoring of greenhouse conditions, as well as control of water supply, humidity, and supplemental lighting. However, the system showed limitations in its ability to regulate higher temperatures during hotter days. For future development of an automated greenhouse in the warmer climate of South Africa, alternative cooling strategies could be explored, such as motorized shading or the integration of an evaporative cooling pad. In addition, the sizing of intake and exhaust fans could be optimized with airflow capacity in mind, favoring an overcompensated design to ensure adequate ventilation under peak heat conditions.

\section*{Conclusions}

The design proved effective in monitoring all key parameters required to sustain healthy plant growth. Automated heating successfully maintained temperatures above the setpoint during cooler nights, while the misting system provided sufficient control of both humidity and water supply. 

The vent and intake fans contributed to improved airflow, helping to limit temperature rise compared to conditions without these actuators. Although overheating remained a challenge, the greenhouse system demonstrated the ability to regulate a microclimate capable of supporting plant life, while providing the user with real-time and historical data and control through the wireless dashboards.

In addition, the solar power system enabled fully off-grid operation, allowing the greenhouse to run entirely on renewable energy. This not only reduced the system’s carbon footprint but also highlighted its potential as a sustainable solution to address food insecurity by providing controlled growing conditions, with the primary limitation being high-temperature management.







